%% ScienceEcosystem — Future Vision
%% Section: The Living Science Record
%% For inclusion in the main business plan document

% ─────────────────────────────────────────────────────────────────────────────
\section*{The Future of ScienceEcosystem: Toward a Living Science Record}
% ─────────────────────────────────────────────────────────────────────────────

\subsection*{The Problem We Are Actually Solving}

Science does not have a memory problem. It has a conversation problem.

Every year hundreds of thousands of papers are published. Most of them respond
to, correct, challenge, or extend earlier work. But the paper that started the
conversation sits frozen — a PDF issued on a fixed date, unchanged, carrying no
trace of what happened to its claims after publication.

A reader encountering that paper today has no way of knowing whether its central
finding was later replicated or failed replication; whether a key method was
subsequently shown to be flawed; whether three papers challenged its
interpretation; whether the authors themselves issued a correction; or whether
the field has quietly moved on. This information exists. It is scattered across
citing papers, letters to editors, conference presentations, informal online
discussions, and private knowledge held by researchers who have followed the
field for years. It is simply never assembled in one place, attached to the
original claim.

The result is that the scientific record looks cumulative but is not navigable
as one. Each paper is an island. The ocean between them — the actual scientific
discourse — is invisible.

This is not a technical failure. It is a structural one. The infrastructure of
science — journals, citation indices, institutional repositories — was designed
to preserve and distribute individual publications, not to represent the
evolving state of knowledge. It does the former well and the latter not at all.

\subsection*{The Current System's Three Structural Failures}

\textbf{1. Publication bias.}
Positive results are published; null results and failures to replicate largely
are not. A paper claiming an effect enters the record permanently. The five
papers that subsequently failed to reproduce it are either unpublished or
scattered, with no mechanism to attach their findings to the original claim.
The record accumulates a systematic distortion that individual researchers must
correct privately, by reading widely and remembering selectively.

\textbf{2. Fragmented discourse.}
Scientific arguments unfold across dozens of papers, over years, in multiple
journals. There is no place where the argument is assembled and made visible as
an argument. A junior researcher entering a field, a policymaker trying to
understand the evidence, or a journalist reporting on a contested finding must
reconstruct the debate manually — if they have the time and the access. Most
do not.

\textbf{3. Static authority.}
Once a paper is published it cannot be updated. Corrections are published as
separate documents, rarely read by those who read the original. Retractions are
indexed in a separate database that most literature searches do not surface.
The scientific record accumulates errors that are technically resolved but
practically invisible to all but the most careful readers.

All three failures share the same root cause: papers are treated as terminal
documents rather than contributions to an ongoing conversation. The publication
event is the end of the process. ScienceEcosystem treats it as the beginning.

\subsection*{The Vision: Science as a Living Record}

The long-term goal of ScienceEcosystem is to build the connective layer that
currently does not exist: a structured, navigable, continuously updated record
of what the scientific community has claimed, tested, challenged, and come to
believe — organised not just by paper, author, and journal, but by
\emph{claim}.

\subsubsection*{The Claim as the Unit of Science}

The fundamental unit of science is not the paper. It is the claim.

A paper makes a number of specific assertions: that a compound inhibits a
pathway at a given concentration; that one method outperforms another on a
specific task; that a correlation holds in a particular population. These claims
have independent fates. Some accumulate strong support. Some are contradicted by
subsequent work. Some are refined and qualified. Some are never tested again.

The scientific record as it currently exists does not track these fates. A
citation index knows that Paper B cited Paper A. It does not know whether Paper
B supported, challenged, or corrected the specific claim in Paper A that it was
responding to. A citation count measures attention, not the direction or quality
of the scientific response.

ScienceEcosystem's long-term architecture treats claims as first-class objects
in the knowledge graph. Each claim is:

\begin{itemize}
  \item Linked to the paper and section where it was made
  \item Linked to every subsequent paper that tested, challenged, supported,
        or qualified it, with the nature of that relationship recorded
  \item Annotated with the current state of evidence — supported, contested,
        refuted, replicated, qualified — based on the accumulated weight of
        subsequent work
  \item Readable by anyone; editable by verified researchers under defined
        governance rules
  \item Traceable: every annotation cites the source it derives from
\end{itemize}

This is not a replacement for papers. Papers remain the authoritative record of
methodology, context, and detail. The claim layer sits above them — a
navigable, living map of what the community currently believes and on what
evidence it rests.

\subsubsection*{The Topic Page as Field Record}

ScienceEcosystem's topic pages are already the closest existing approximation
of this vision. They currently aggregate a Wikipedia overview, the most
associated papers, key authors, and publication and citation trends. They are
already more useful than a raw literature search for understanding a field.

The mature version of a topic page becomes the living record of a scientific
field. It shows:

\begin{itemize}
  \item The central claims and their current evidential status
  \item Where the field has reached consensus and where active debate continues
  \item The methodological fault lines — which results depend on which methods,
        and what it means when methods produce different outcomes
  \item The open questions: what has been called for but not yet investigated
  \item A timeline of how understanding evolved — what was believed in 2010,
        what changed and why, what the current picture is
\end{itemize}

This is the knowledge that currently exists only in the heads of domain experts.
A researcher who has spent five years in a field holds a detailed mental model
of the terrain — the reliable findings, the contested claims, the known
artefacts of specific methods. That model is never written down in a single
place. The senior researcher retires; the model is lost. ScienceEcosystem makes
it structural, citable, and updatable.

\subsubsection*{The Living Paper Format}

The culmination of this vision is a publication format whose native structure
is the knowledge graph rather than the linear document. A researcher submitting
a paper to ScienceEcosystem writes prose, adds figures, uploads data, and runs
code — the experience is that of writing a paper. Behind the surface, every
entity, method, result, claim, and provenance relationship is stored as a node
or edge in the graph.

The rendered output — the paper that readers see — is generated from the graph
and can be rendered as a traditional article, a structured dataset, a summary,
or an AI-queryable object. The graph is the truth; the paper is a view of it.

This means that when a subsequent paper challenges one finding in the original,
that challenge attaches to the specific claim node, not to the paper as a whole.
The paper does not become invalid; one of its claims is updated in the record.
The conversation is visible, structured, and permanently attached to the science
it concerns.

\subsection*{The Strategic Approach: Incremental and Useful at Every Stage}

The central strategic insight of ScienceEcosystem's roadmap is that the path to
transforming scientific infrastructure is to be genuinely useful at every
intermediate stage — not to build the full vision first and then ask researchers
to adopt it.

Most projects attempting to reform scientific publishing fail for the same
reason: they build the ideal system first, then discover that researchers will
not change their workflow until the system is already the norm. They need the
community to adopt the tool, but the community will not adopt the tool until
the community has already adopted the tool.

ScienceEcosystem avoids this by starting with what researchers already need —
a better way to find, read, and track papers — and expanding from there. Each
step is an improvement on what exists. The infrastructure for the next step is
built by building the current step well. By the time the platform reaches the
full claim-layer vision, researchers are not being asked to change how they
work; they are being offered an extension of something they already depend on.

The four stages of this progression are described below.

\subsubsection*{Stage 1 — Useful Today}

A better place to find, read, and track papers than any existing alternative.
Clean author profiles, rich topic pages, a personal library, and a browser
extension that connects any publisher page or PDF to the platform's data.
This stage builds the user base, establishes credibility, and demonstrates
that the platform is maintained and reliable. No vision is required at this
stage — it simply has to work better than what exists.

\textbf{Status:} Substantially complete. The platform is live with all core
features functional.

\subsubsection*{Stage 2 — Making the Conversation Visible}

The first step that does not exist anywhere else. Citation context and stance
classification: for each paper, show the actual sentences in which it is cited
by other papers, and classify the relationship as supporting, challenging,
correcting, extending, methodological, or background. Add an author response
field — the cited author can attach a note visible to all readers, creating
the first mechanism for direct scientific dialogue on the publication record.

This stage produces the first layer of claim-level data. It requires no
structural change to how research is conducted or published. It works
retroactively on the existing literature. And it gives researchers something
they genuinely want: not just a count of citations, but a sense of how their
work has been received.

\textbf{Technical foundation:} Semantic Scholar API (citation context
sentences); LLM-based stance classification (short, cheap inference on
sentence-level input); author response stored per paper in the platform
database.

\textbf{Status:} Planned. Buildable with current infrastructure and no
additional funding.

\subsubsection*{Stage 3 — Topic Pages as Field Records}

Once stance-classified citation data accumulates at scale, it can be
aggregated at the topic level. A topic page begins to show not just the
papers and authors in a field but the current state of its arguments —
where evidence is strong, where it is contested, where claims have been
made but not adequately tested. Open questions become visible as a category,
not just as a sentence at the end of individual papers.

This stage requires no new data sources — it is an aggregation and
presentation layer over the citation stance data built in Stage 2. It is
however where the product becomes qualitatively different from everything
that exists: the first place where a non-expert can genuinely understand the
evidential state of a scientific field without reading hundreds of papers.

\textbf{Status:} Dependent on Stage 2 data accumulation. Medium-term target.
Requires modest infrastructure investment and possibly a small team.

\subsubsection*{Stage 4 — Claims as First-Class Objects}

The full realisation of the vision. Claims extracted from the literature —
beginning with abstracts and results sections, using LLM-based extraction —
become nodes in the knowledge graph. Claims from different papers that address
the same assertion are linked. The graph represents not papers and citations
but assertions and their evidential histories.

Topic pages at this stage show field-level consensus maps. The living paper
format becomes available for new submissions. The platform becomes a genuine
alternative publication venue for structured science, not just an indexing
layer above existing journals.

This stage requires significant infrastructure investment, a functioning team,
and a community of researchers willing to engage with the annotation and
governance systems. It is a long-term target, not a near-term deliverable.

\textbf{Status:} Long-term vision. Requires funding and team. The earlier
stages build toward this in both technical infrastructure and community trust.

\subsection*{The Knowledge Graph Architecture}

The knowledge graph underlying this vision has multiple node types and edge
types that accumulate progressively across the four stages.

\textbf{Existing nodes (Stage 1):}
Paper $\cdot$ Author $\cdot$ Topic $\cdot$ Institution $\cdot$ Journal

\textbf{Added in Stage 2:}
CitationContext $\cdot$ Stance $\cdot$ AuthorResponse

\textbf{Added in Stage 3:}
FieldConsensus $\cdot$ OpenQuestion $\cdot$ MethodologicalLineage

\textbf{Added in Stage 4:}
Claim $\cdot$ Method $\cdot$ Dataset $\cdot$ EvidenceStatus $\cdot$ ReplicationRecord

\textbf{Key edges:}
\texttt{cites\_with\_stance(Paper, Paper, Stance)} $\cdot$
\texttt{makes\_claim(Paper, Claim)} $\cdot$
\texttt{tests\_claim(Paper, Claim, Result)} $\cdot$
\texttt{uses\_method(Paper, Method)} $\cdot$
\texttt{produces\_dataset(Paper, Dataset)} $\cdot$
\texttt{responds\_to(AuthorResponse, CitationContext)}

The graph is intentionally built on standard relational database infrastructure
(PostgreSQL) in the early stages, with graph traversal handled through
recursive queries. This avoids premature investment in specialised graph
database infrastructure while maintaining the structural integrity needed for
the later stages.

\subsection*{On AI and the Knowledge Graph}

A common question in building this kind of system is whether a knowledge
graph is needed at all, or whether large language models can simply answer
queries about the scientific literature on demand.

The answer is that both are needed, and they serve different functions.

LLMs are used for extraction: reading citation context sentences and
classifying stance, extracting claims from abstracts, identifying
methodological relationships. They are fast, cheap for short inputs, and
continuously improving. They are the tool that makes it practical to process
the existing literature at scale.

The knowledge graph is needed for storage, provenance, and queryability. An
LLM cannot tell you reliably what the current state of evidence on a specific
claim is across ten thousand papers — it has no ground truth, only training
data. The graph stores the results of extractions with their sources, allows
structured queries across the full literature, and provides the auditable
record that scientific infrastructure requires.

The combination — AI for extraction, graph for storage and traversal — is the
architecture that makes the vision achievable. Neither alone is sufficient.

\subsection*{Why This is Different From Existing Approaches}

Several projects have attempted parts of this vision. None has built the full
connective layer.

\textbf{Semantic Scholar} classifies citation intent (background, methodology,
result) but not stance (agrees, challenges, corrects). It does not expose
claims or support author response. It is a research tool, not a community
infrastructure.

\textbf{PubPeer} supports post-publication annotation but is structured as a
comment thread attached to a paper identifier. It has no stance classification,
no claim linking, no topic-level aggregation, and no governance model for
claim status.

\textbf{ResearchGate} is a social network with paper hosting. Its primary
value is access to paywalled PDFs. It has no claim layer, no citation stance,
and its incentive model (gamified RG Score) does not align with surfacing
contested or failed results.

\textbf{Wikipedia} aggregates field knowledge but is not structured at the
claim level, does not link to primary literature systematically, has no
mechanism for representing contested claims with their evidential basis, and
is written by volunteers rather than the researchers who produced the work.

ScienceEcosystem occupies a position none of these occupy: a platform that
starts with the existing literature and the researchers who produced it, and
builds the connective layer — claim by claim, stance by stance — that makes
the scientific discourse navigable.

\subsection*{The Infrastructure Perspective}

It is worth stating directly what kind of project ScienceEcosystem is.

It is not a product in the commercial sense. It is infrastructure — like a
library, a road, or a communication protocol. Its value is not in what it sells
but in what it enables. Infrastructure of this kind has specific properties:
it is more valuable when it is open (because openness maximises the number of
people who can build on it), it should not be owned by any single institution
(because dependence on a single owner creates fragility and potential for
capture), and it should be designed to outlast the organisations and
technologies of the moment it was built in.

This is why the non-profit structure, the open data principles, and the
governance model described elsewhere in this plan are not secondary to the
technical vision — they are constitutive of it. An infrastructure that can be
acquired, paywalled, or shut down when it becomes inconvenient is not
infrastructure. It is a product with a temporary public benefit.

The living science record, as described here, is only possible if researchers
and institutions trust that it will be there in fifty years, that the data they
contribute will not be sold or withheld, and that the governance of what counts
as good evidence will not be controlled by those with a commercial interest
in the answer.

That trust is built slowly, step by step, in exactly the way this roadmap
proposes: by being reliable before being ambitious, by being useful before
asking for commitment, and by demonstrating at every stage that the platform
serves science rather than extracting value from it.

\section*{3. Memory Alpha}
Scientific knowledge is often lost not because data disappears, but because context, 
methods, and meaning become disconnected over time. When datasets, code, 
and publications are stored separately, research becomes difficult to interpret, 
reuse, or reconstruct as technologies and institutions change.

ScienceEcosystem addresses this by acting as a connective knowledge layer. 
It links data, code, analyses, and publications into coherent, 
self-describing research objects that preserve provenance and assumptions alongside results.

By prioritising meaning over storage, ScienceEcosystem creates a foundation that allows 
scientific knowledge to remain understandable and resilient in the long term. 
As this connective layer matures, it can evolve into a broader scientific memory infrastructure, 
supporting future preservation and continuity initiatives without being tied to any single 
technology or institution.

